\documentclass[10pt,a4paper]{article}
\pdfmapfile{+lm.map}
\pdfmapfile{+cm.map}
\pdfmapfile{+symbols.map}
\usepackage[T1]{fontenc}
\usepackage[utf8]{inputenc}
\usepackage[ngerman,provide=*]{babel}
\usepackage{lmodern,amsmath,amssymb,amsthm,booktabs,array,geometry,xcolor,hyperref,fancyhdr,microtype}
\geometry{left=23mm,right=23mm,top=24mm,bottom=24mm}
\definecolor{navy}{HTML}{17334B}
\definecolor{teal}{HTML}{087E8B}
\hypersetup{colorlinks=true,urlcolor=teal,linkcolor=navy,pdftitle={CST-DUST: Gemeinsamer minimaler Skalarwirkungsansatz v0.1},pdfauthor={X6 Forschungsprogramm / Arbeitsausarbeitung}}
\pagestyle{fancy}\fancyhf{}
\fancyhead[L]{\small\sffamily\color{navy}CST-DUST / Gemeinsamer Skalarwirkungsansatz}
\fancyhead[R]{\small\sffamily v0.1}
\fancyfoot[L]{\footnotesize 11. November 2026 / Physikalische Qualifikation: HOLD}
\fancyfoot[R]{\thepage}
\setlength{\headheight}{14pt}
\setlength{\parindent}{0pt}\setlength{\parskip}{6pt}
\setlength{\emergencystretch}{2em}
\newcommand{\R}{\mathbb R}
\newcommand{\dd}{\mathrm d}
\newcommand{\ii}{\mathrm i}
\newcommand{\sgn}{\operatorname{sgn}}
\newcommand{\sdet}{\operatorname{sqrt}_{\!c}\det}
\newcommand{\status}[1]{\textsf{\textbf{#1}}}
\newcommand{\sect}[1]{\section*{\color{navy}#1}}
\newcommand{\smallurl}[1]{{\footnotesize\url{#1}}}
\begin{document}
{\sffamily\small\color{teal}SCHRITT 1 / ANSCHLUSS AN PR 278}
\vspace{5mm}

{\sffamily\bfseries\color{navy}\LARGE Gemeinsamer minimaler\par Skalarwirkungsansatz\par}
\vspace{3mm}
{\large CST-DUST-KG-MIN-v0.1: definierter Kandidat, bedingter Grenzfall\par}

\sect{Ergebnis und physikalische Rolle}
Schritt~1 legt einen konkreten gemeinsamen \emph{Wirkungskandidaten} für
Geometrie, wirbelfreien Staub und ein zusätzliches reelles Skalarfeld fest.
Das Feld $\Phi$ beschreibt eigenständige Materie. Es ist weder die Staubuhr
$T$ noch eine bereits hergeleitete freie Anregung dieser Uhr.
Die Motivation ist das bekannte klassische Modell aus Einstein-Gravitation,
wirbelfreiem Staub und einem minimal gekoppelten Klein-Gordon-Feld.

Für $v=\ell^4>0$, $\mu=m/\hbar\geq0$ bei $c=1$ und die unten definierte
diskrete Produktform $\Gamma$ lautet der Kandidat
\[
\boxed{\begin{aligned}
S_{\rm tot}={}&S_g(C)+S_d(C,T,w)+S_\Phi(C,\Phi)+S^g_{\rm outer},\\
S_d={}&-\frac12\sum_xw_x\bigl(\Gamma_x[T]+1\bigr),\\
S_\Phi={}&-\frac12\sum_xv\bigl(\Gamma_x[\Phi]+\mu^2\Phi_x^2\bigr).
\end{aligned}}
\]
Die zusätzliche Materie koppelt an dieselbe Geometrie. Eine direkte
$T,w,\Phi$-Wechselwirkung wird in diesem minimalen Kandidaten auf null
gesetzt; dies ist eine Modellwahl. Die physikalische äußere
Gravitationsrandwirkung bleibt offen. Ihr bisheriger Kontrollwert null
wird nicht als hergeleitetes Randgesetz ausgegeben.

\begin{center}
\begin{tabular}{p{.51\linewidth}p{.38\linewidth}}
\toprule
Ergebnis & Status\\\midrule
Endliche gemeinsame Wirkungsform & ausdrücklich definiert\\
Symmetrische Feld-Hesse-Matrix & aus der Wirkung hergeleitet\\
Skalare Randreaktion / Summationsdefekt & exakt bestimmt\\
Klassische Zielgleichungen & analytisch geprüft\\
33 angeschlossene $AA$-Kontrollhistorien & Feldmatrizen exakt exportiert\\
Automatisierte Kontrollen & 25 Tests bestanden\\
Physikalische $8\times8$-Materiematrix & \texttt{NOT\_CONSTRUCTED}\\
Alle acht physikalischen Gates & \status{HOLD}\\\bottomrule
\end{tabular}
\end{center}

\textbf{Was nicht behauptet wird.} Dieser Ansatz ist durch Feldinhalt,
Symmetrien und einen klassischen Zielgrenzfall motiviert, aber nicht
eindeutig aus CST abgeleitet. Konvergenz der gemeinsamen Wirkung und ihrer
Variationen, physikalische Zustände, Quantenmaß, Stabilität und äußere
Gravitationsrandphysik sind nicht bewiesen. Eine zusätzliche Materiespezies
ist eine Erweiterung des Staubmodells und keine Erklärung des Staubs selbst.

Das Paket berechnet keine neue physikalische Übergangsamplitude. Es ersetzt
die bislang beliebige symbolische Eingabe $H_C(T,w)$ durch eine überprüfbare,
explizit deklarierte Kandidatenmatrix $H_C$.

\newpage
\sect{1. Feldinhalt, Einheiten und bewusst gewählte Annahmen}
Wir verwenden die Signatur $(-+++)$ und $c=1$. Mit $[S]$ als
Wirkungsdimension gelten
\[
[T]=L,\quad [w]=[S],\quad [v]=L^4,\quad [B]=L^{-2},
\quad [\Phi]=[S]^{1/2}L^{-1},\quad[\mu]=L^{-1}.
\]
$w=v\rho$ ist das geerbte integrierte Staubdichtegewicht. $v$ ist ein
geometrisches Volumengewicht und bleibt auch bei $w=0$ positiv.
Jeder Term in $S_d$ und $S_\Phi$ besitzt die Dimension Wirkung.
Insbesondere wird die Skalarwirkung nicht mit $w$ anstelle von $v$ gewichtet.
Der Dichte-Nullzweig bezeichnet Abwesenheit von Staub; er bedeutet nicht
Abwesenheit des zusätzlichen Skalarfelds oder dessen Vakuumzustand.

\textbf{Warum eine Staubuhrstörung kein freies KG-Feld ist.}
Betrachte allein die geerbte \emph{klassische Kontinuumswirkung} auf festem
Minkowski-Hintergrund. Setze $T=t+\varepsilon\pi$ und
$\rho=\rho_0+\varepsilon\sigma$, mit konstanter $\rho_0>0$. Bis zur
zweiten Ordnung ist die Lagrangedichte
\[
\mathcal L_d=
\varepsilon\rho_0\dot\pi+
\varepsilon^2\left[\frac{\rho_0}{2}
\bigl(\dot\pi^2-|\nabla\pi|^2\bigr)+\sigma\dot\pi\right]
+O(\varepsilon^3).
\]
Die erste Ordnung ist unter den genannten Voraussetzungen eine totale
Zeitableitung. Die Variation des quadratischen Anteils nach $\sigma$
erzwingt $\dot\pi=0$. Die scheinbar freie Wellengleichung nach Einfrieren
von $\sigma$ wäre daher nicht die gemeinsame Staubvariation.
Diese einfache Kontrolle begründet die Trennung von $\Phi$ und $T$;
sie ist kein vollständiger Staub-Störungsformalismus mit dynamischer Metrik.

\textbf{Minimaler Kandidat.} Gewählt werden ein reelles Feld, eine
quadratische Wirkung, $\mu^2\geq0$, minimale Krümmungskopplung $\xi=0$,
keine direkte Staub-Skalar-Kopplung und keine Skalar-Selbstwechselwirkung.
Damit bleiben Gesamtfeldparität $\Phi\mapsto-\Phi$ und die geerbte
Uhrverschiebung $T\mapsto T+c$ auf Wirkungsebene erhalten. Bei $\mu=0$
ist zusätzlich $\Phi\mapsto\Phi+c_\Phi$ eine exakte Wirkungssymmetrie.
Randzustände und Regulatoren müssen diese Symmetrien gegebenenfalls separat
respektieren. Weder Parität noch Kovarianz erzwingen $\xi=0$ oder das
Fehlen von Selbstwechselwirkungen.

Ein beispielsweise hinzugefügter Term
$-\frac g{2F^2}\sum_xw_x\Phi_x^2$, mit $[F]=[\Phi]$, würde die positive
Dichtevariation zu
\[
\frac{\partial S}{\partial w_x}
=-\frac12\left(\Gamma_x[T]+1+\frac{g\Phi_x^2}{F^2}\right)
\]
ändern. Das wäre ein anderes Modell mit veränderter Staubbedingung.
Ein solcher Term wird nicht zur Erzeugung eines gewünschten Mischblocks
eingeführt. Auch die Werte von $\ell$ und $\mu$ werden hier nicht angepasst.

Die getrennte Verwendung von wirbelfreiem Staub und zusätzlicher Materie
besitzt einen bekannten Kontinuumsrahmen [L1]. Dessen kanonische oder
LQG-Quantisierung wird dadurch nicht auf unser CST-Paket übertragen.

\newpage
\sect{2. Diskrete gemeinsame Wirkung und exakter Feld-Hessian}
Auf derselben endlichen strikten Ordnung $C$ wie bisher definiere
\[
a=\frac4{\sqrt6\ell^2},\quad B_{xx}=-a,\quad
B_{xy}=a\,b\bigl(k_C(y,x)\bigr)\ (y\prec x),
\quad b=(1,-9,16,-8,0,\ldots).
\]
Es wird der geerbte unverschmierte 4D-BD-Kandidat benutzt. Diese
endliche Festlegung belegt für sich keine 4D-Manifoldlikeness der kleinen
Kontrollordnungen. Für Felder $f,h$ setze die polarisierte Produktform
\[
\Gamma_x[f,h]=\frac12\left[B(fh)_x-f_x(Bh)_x-h_x(Bf)_x
+f_xh_x(B1)_x\right].
\]
Alle Produkte in dieser Zeile sind punktweise. Direktes Ausmultiplizieren
liefert die exakte, diagonalenunabhängige Identität
\[
\boxed{\Gamma_x[f,h]=\frac12\sum_{y\prec x}
B_{xy}(f_x-f_y)(h_x-h_y).}
\]
Wir schreiben $\Gamma[f]=\Gamma[f,f]$. Daraus folgt
\[
S_\Phi=-\frac14\sum_{y\prec x}v_xB_{xy}(\Phi_x-\Phi_y)^2
-\frac{\mu^2}{2}\sum_xv_x\Phi_x^2.
\]
Die algebraischen Formeln gelten auch für vorgegebene positive $v_x$;
das konkrete Modell verwendet durchgängig $v_x=v=\ell^4$.

Sei $D_{\rm inc}$ die Inzidenzmatrix \emph{aller berücksichtigten kausalen
Paare}, mit $-1$ am Vorgänger und $+1$ am Nachfolger, und
$W_E=\operatorname{diag}(v_xB_{xy})$, $V=\operatorname{diag}(v_x)$.
Dann gilt exakt
\[
\boxed{S_\Phi=\frac12\Phi^TH_C\Phi,
\qquad H_C=-\frac12D_{\rm inc}^TW_ED_{\rm inc}-\mu^2V.}
\]
$H_C$ ist reell-symmetrisch, obwohl $B$ retardiert ist. Das folgt aus
der Variation der definierten Wirkung und nicht aus einer nachträglich
auferlegten Symmetrisierung einer Bewegungsgleichung.
$\nabla_\Phi S_\Phi=H_C\Phi$ und $H_C1=-\mu^2v$ folgen unmittelbar.
Die wechselnden BD-Gewichte machen diese Raumzeitmatrix im Allgemeinen
indefinit; sie ist kein positiver räumlicher Energieoperator.

\textbf{Gemeinsame Gravitation.} Bei konstantem $v$ gilt
$\sum_x(B1)_x=-aF_4(C)$. Mit $\widehat R_x=-2(B1)_x$ ist daher
\[
\frac1{16\pi G}\sum_xv\widehat R_x
=\frac{\ell^2}{2\pi G\sqrt6}F_4(C)=\hbar\kappa F_4(C)=S_g(C),
\]
mit der geerbten Konvention $\ell_P^2=8\pi G\hbar$.
Das ist eine endliche Normierungsidentität, kein Nachweis des vollständigen
Einstein-Hilbert-Grenzwerts samt äußerem Rand. Die korrigierte Produktform
verhindert zugleich eine unbeabsichtigte $R\Phi^2$-Kopplung im Zielgrenzfall.

\newpage
\sect{3. Randkorrektur und das Variationsproblem}
Die bekannte Behauptung $S_\Phi=\frac12\Phi^TVB\Phi$ würde sowohl
Nullordnungsterm als auch Randsummation ungeprüft übernehmen. Die genaue
Beziehung folgt aus
\[
r=B1,\qquad D_C=B-\operatorname{diag}(r),\qquad
\tau=B^Tv-Vr=D_C^Tv.
\]
Insbesondere gilt $D_C1=0$ und für jede Testfunktion $f$
$\sum_xv_x(D_Cf)_x=\tau^Tf$. Durch Einsetzen in die Produktform erhält man
\[
\sum_xv_x\Gamma_x[\Phi]=
\frac12\tau^T\Phi^{\circ2}-\Phi^TVD_C\Phi,
\]
also
\[
\boxed{S_\Phi=
\frac12\Phi^TV(D_C-\mu^2I)\Phi
-\frac14\tau^T\Phi^{\circ2}.}
\]
Dabei bezeichnet $\Phi^{\circ2}$ das komponentenweise Quadrat.
Der zugehörige Hessian lautet äquivalent
\[
H_C=\frac12(VD_C+D_C^TV)
-\frac12\operatorname{diag}(\tau)-\mu^2V.
\]
Der Summationsdefekt $\tau$ ist exakt bestimmt. Bei endlicher nichtlokaler
Ordnung kann er über viele Elemente verteilt sein. Er wird nicht ohne
Grenzwertnachweis als lokaler GHY-Term oder als ausschließlich auf dem
äußeren Rand getragener Vektor bezeichnet.

\textbf{Kontinuumsvergleich des Terms.} Falls der integrierte Operator
$D_C$ zu $\Box$ konvergiert und die Randintegration kontrolliert ist,
entspricht $-\tau^T\Phi^{\circ2}/4$ dem Ausdruck
\[
-\frac14\int_{\partial M}\dd\Sigma_a\,\nabla^a(\Phi^2)
=-\frac12\int_{\partial M}\dd\Sigma_a\,\Phi\nabla^a\Phi.
\]
Damit ergibt sich die gewöhnliche partielle Integration der skalaren
Gradientenwirkung. $\dd\Sigma_a$ bezeichnet das orientierte Flächenelement;
für nullartige Ränder wird keine Einheitsnormalen-Konvention unterstellt.

\textbf{Endliche Randvariation.} Für eine ausdrücklich festgelegte Zerlegung
$(b,q)$ in äußere und innere Feldwerte schreibe
$H=\bigl(\begin{smallmatrix}B_\partial&J^T\\J&A\end{smallmatrix}\bigr)$.
Dann ist
\[
\delta S_\Phi=(B_\partial b+J^Tq)^T\delta b+(Jb+Aq)^T\delta q.
\]
Bei festen äußeren Werten gilt $\delta b=0$; innere Stationarität verlangt
$Aq+Jb=0$. Die äußere Reaktion $p_b=B_\partial b+J^Tq$ wird beibehalten.
Es ist für diese skalare Dirichlet-Variation kein weiterer Randterm nötig.
Bei Neumann-Daten, Paketprojektion oder einer anderen Randdefinition ist
ein eigener Vertrag erforderlich. Die Eignung des äußeren Kragens als
vollständige physikalische Randinformation gehört zu Schritt~2.

Die offene \emph{Gravitations}randwirkung $S^g_{\rm outer}$ wird durch
diese skalare Rechnung nicht bestimmt. Ein zusätzlicher skalarer Term
$-\tau^T\Phi^{\circ2}/4$ darf zur bereits vollständigen Paarwirkung nicht
nochmals addiert werden.

\newpage
\sect{4. Klassischer Zielgrenzfall und seine Voraussetzungen}
Der Zielgrenzfall ist die gemeinsame Kontinuumswirkung
\[
S_{\rm cl}=\frac1{16\pi G}\int_M\sqrt{-g}\,R\,\dd^4x
-\frac12\int_M\sqrt{-g}\,
\left\{\rho\bigl[(\nabla T)^2+1\bigr]
+(\nabla\Phi)^2+\mu^2\Phi^2\right\}\dd^4x+S^g_{\partial M}.
\]
Der wirbelfreie Staubterm und unabhängige zusätzliche Materie sind ein
bekannter klassischer Ausgangspunkt [3]. Der vorliegende diskrete Ansatz
ist dadurch motiviert; seine gemeinsame CST-Konvergenz folgt daraus nicht.

\textbf{Produktform.} Angenommen, $B$ hat für die benötigten glatten
Testfelder den Grenzoperator $\Box-R/2$. Dann ergibt Einsetzen in die
\emph{polarisierte} Definition von $\Gamma$
\[
\Gamma[f,h]\longrightarrow\tfrac12\{
(\Box-R/2)(fh)-f(\Box-R/2)h
-h(\Box-R/2)f+fh(-R/2)\}
=\nabla_af\nabla^ah.
\]
Die Krümmungsanteile heben sich auf. Insbesondere entsteht hieraus keine
vorgegebene nichtminimale $R\Phi^2$-Kopplung.
Die Literatur beschreibt den entsprechenden \emph{mittleren}
Operatorgrenzfall unter ihren Voraussetzungen [1,2]. Das ist allein kein
Satz über die gemeinsame Zufallswirkung oder deren erste Variation.

\textbf{Schwache skalare Variation.} Für eine Testvariation $\eta$ gilt exakt
\[
\delta_\eta S_\Phi=-\sum_xv\,\Gamma_x[\Phi,\eta]
-\mu^2\sum_xv\Phi_x\eta_x.
\]
Erforderlich sind Konvergenz dieser integrierten polarisierten Form,
Volumenquadratur und eine kontrollierte Randbehandlung. Erst dann folgt
für kompakt getragene $\eta$ die schwache Zielgleichung
\[
\boxed{(\Box-\mu^2)\Phi=0.}
\]
Eine bloße punktweise Mittelwertübereinstimmung ersetzt diese Annahmen
nicht. Auch Varianz- und Fluktuationskontrolle sind noch ausstehend.

\textbf{Staub und Rückwirkung im Zielmodell.} Unabhängige Variation nach
$\rho$ und $T$ liefert
\[
(\nabla T)^2=-1,\qquad \nabla_a(\rho\nabla^aT)=0.
\]
Mit $u_a=-\nabla_aT$ folgt geodätische Staubbewegung, solange das glatte
wirbelfreie Staubmodell gilt, insbesondere vor seinen Kaustiken.
Für zusätzlich kontrollierte Geometrievariationen lautet das Ziel
\[
G_{ab}=8\pi G\left[\rho T_aT_b+\Phi_a\Phi_b
-\tfrac12g_{ab}\bigl((\nabla\Phi)^2+\mu^2\Phi^2\bigr)\right].
\]
Hier werden die Kontinuums-Materiefelder bei der Metrikvariation festgehalten;
die Staubnorm wird danach eingesetzt. Das ist nicht dieselbe Operation wie
eine diskrete Variation bei festem $w=v\rho$ und veränderlichem $v$.
Die Formel ist ein analytisch geprüfter Ziel-Stresstensor, kein Nachweis
diskreter gravitativer Rückwirkung.

Auf festem Minkowski-Hintergrund sind
$\omega^2=|\boldsymbol{k}|^2+\mu^2$ und
$\mathcal E_\Phi=(\dot\Phi^2+|\nabla\Phi|^2+\mu^2\Phi^2)/2\geq0$
geprüfte klassische Zielidentitäten. Sie beweisen keine Stabilität des
endlichen CST-Operators. Der diskutierte ursprüngliche retardierte
Vierdimensional-Operator besitzt zudem bekannte Stabilitätsprobleme [4];
der hier variierte symmetrische Operator benötigt eine eigene Prüfung.

\newpage
\sect{5. Kopplung über Geometrie und Anschluss an die Historien}
Bei festem $C$ separieren Staub und zusätzliches Skalarfeld in der Wirkung.
Unter \emph{zusätzlich separierbaren} Integrationsmaßen und Randpaketen
hat der Beitrag einer Historie deshalb die Form
\[
e^{\ii S_g(C)/\hbar}\,D_C^{\rm dust}\,A_C^\Phi.
\]
Die gemeinsame Summe über Geometrien separiert im Allgemeinen nicht:
\[
\sum_C\nu_Ce^{\ii S_g(C)/\hbar}D_C^{\rm dust}A_C^\Phi
\ne
\left(\sum_C\nu_Ce^{\ii S_g(C)/\hbar}D_C^{\rm dust}\right)
\left(\sum_C\nu_Ce^{\ii S_g(C)/\hbar}A_C^\Phi\right).
\]
Normierungen sind dabei separat festzulegen. Bereits für eine normalisierte
Zwei-Geometrien-Summe mit Gewichten $p,1-p$ ist die Differenz
$p(1-p)(D_1-D_2)(A_1-A_2)$. Diese algebraische Aussage beschreibt die
Möglichkeit geometrisch vermittelter Korrelation; sie ist noch kein
Nachweis eines physikalischen Übergangs oder einer eindeutigen Ursache.

\textbf{Kein künstlicher Entkopplungsparameter.}
Der Kandidat besitzt keine direkte Staub-Skalar-Kopplung. Eine Skalierung
$\alpha S_\Phi$ mit anschließendem $\alpha\to0$ ist kein bereits
qualifizierter Entkopplungsvergleich: Bei unbeschränkten reellen inneren
Feldachsen verschwindet dann die feldabhängige Wirkung und unkontrollierte
Integrationsvolumina bleiben zurück. Auch ein willkürlich gesetzter
$q_{0+}$-Eintrag ist ausgeschlossen. Ein späterer Referenzvertrag muss
etwa die Geometrie festhalten oder einen ausdrücklich statistischen
Vergleich mit unabhängigen Geometrieziehungen definieren; seine
physikalische Interpretation bleibt zu begründen.

\textbf{Gemeinsame endliche Geometrievariation.}
Bei einem zulässigen Ordnungswechsel $C\to C'$ auf derselben Elementmenge,
festem $\ell,v,T,w,\Phi$ und demselben Operatorvertrag ist
\[
\begin{aligned}
\Delta S_{\rm tot}={}&\hbar\kappa\Delta F_4+\Delta S^g_{\rm outer}\\
&-\frac14\sum_{x\ne y}
\left[w_x(T_x-T_y)^2+v(\Phi_x-\Phi_y)^2\right]\Delta B_{xy}.
\end{aligned}
\]
Die Summe verwendet die gemeinsame Matrixindizierung; nicht vorhandene
Relationen haben Gewicht null. Der Massenterm bleibt bei dieser
festen Volumen- und Elementzahlvariation unverändert.
Alle durch den Ordnungswechsel geänderten Intervallgrößen müssen in
$\Delta B$ eingehen, auch für zuvor vorhandene Relationen.
Ein physikalisches Quantengesetz wird nicht durch die Forderung
$\Delta S=0$ für jeden einzelnen diskreten Zug ersetzt.

\textbf{Schnitt und Verklebung.}
Beim Vergleich von $S_\Phi[C]$ mit den Wirkungen zweier induzierter
Teilordnungen entstehen sowohl Paare über den Schnitt als auch Änderungen
der Intervallgewichte innerhalb der Teile. Beide Beiträge werden bereits
von der Paarwirkung erfasst. Eine zusätzliche frei gewählte skalare
Schnittwirkung würde diese Beiträge ohne neuen Vertrag doppelt zählen.
Das ersetzt keine Herleitung des gravitativen Junction-Gesetzes.

Der bereits publizierte Basisvertrag bleibt unverändert:
Strata zuerst, Zeilen ausgehend, Spalten eingehend; die Carrier-Bezeichnung
nennt dagegen zuerst den Eingang. Die skalare Wirkung gilt sowohl im
$w=0$- als auch im $w>0$-Zweig. Sie bestimmt weder deren physikalisches
Maß noch Projektoren, Pakete oder die Mischungsamplitude $K_{0+}$.

\newpage
\sect{6. Ein überprüfbares Resultat für die 33 AA-Kontrollen}
Verwendet wird ausschließlich die bereits deklarierte kleine
Kontrollfamilie mit Eingang $L=\{0,1\}$, Innerem $I=\{2\}$ und Ausgang
$R=\{3,4\}$. Beide äußeren Träger sind Antiketten ($A\to A$).
In der abgedeckten Familie besitzt das innere Element mindestens einen
Vorgänger aus $L$ und einen Nachfolger aus $R$. Setze ihre Anzahlen
$p,q\in\{1,2\}$. Alle an $2$ anliegenden Intervalle sind leer, also
$B_{2y}=a$ beziehungsweise $B_{x2}=a$. Daher folgt direkt aus der Hesse-Matrix
\[
\boxed{H_{II}=-v\left[\frac a2(p+q)+\mu^2\right],\qquad
|H_{II}|\geq v(a+\mu^2)>0.}
\]
Für jedes feste $\ell>0$ und $\mu^2\geq0$ ist der innere Block damit
invertierbar. Ein stationärer innerer Feldwert und der formale Schurblock
sind für diese Kontrollen eindeutig:
\[
q_*=-H_{II}^{-1}Jb,\qquad
H_\partial^{\rm eff}=B_\partial-J^TH_{II}^{-1}J.
\]
Das ist kein bereits normalisierter Quantenkern. Insbesondere sind
Randzustände, Maß und Regulatorgrenzwert dadurch nicht bestimmt.

\begin{center}
\begin{tabular}{lrr}
\toprule
$H_{II}/(va)$, mit $u=\mu^2/a$ & abgedeckte AA & schwache AA\\\midrule
$-1-u$ & 20 & 32\\
$-3/2-u$ & 12 & 16\\
$-2-u$ & 1 & 1\\\midrule
Gesamt & 33 & 49\\\bottomrule
\end{tabular}
\end{center}
Alle 33 Hesse-Matrizen sind exakt rational für $\mu=0$ exportiert;
beliebige nichtnegative Masse wird durch
$H=va(H_0-u\mathbf1_5)$ ergänzt. Die natürliche Elementreihenfolge ist
$(0,1,2,3,4)$; der Gaußadapter verwendet $(0,1,3,4,2)$, also äußere Werte
vor dem inneren Wert. Das ist keine Änderung der achtdimensionalen
Strata-/Randzustandsbasis.

\textbf{Grenzen des Resultats.}
Die Schranke ist nicht gleichmäßig für $\ell\to0$, denn
$va=4\ell^2/\sqrt6\to0$. Sie beweist keine Stabilität bei Verfeinerung.
Außerdem gilt sie nicht für beliebige Kettenränder: Die Fünferkette hat
$H_{II}/(va)=8-u$, also bei $u=8$ einen singulären inneren Block und
gleichzeitig $J\ne0$. Ein einfacher Schur-Inversenansatz ist dort
unzulässig. Dies allein beweist keine Divergenz eines regulierten,
projizierten Gesamtkerns; dafür ist der vorhandene gemeinsame
Nullraumtest mit dem konkreten Projektionsvertrag maßgeblich.

\textbf{Keine Eindeutigkeit der Wirkung.}
Auch bei denselben führenden glatten Zielgleichungen sind Zusatzterme
möglich. Beispielsweise respektiert
\[
\delta S=-\frac{\zeta\ell^2}{2}\sum_xv(D_C\Phi)_x^2,
\qquad
\delta H=-\zeta\ell^2D_C^TVD_C
\]
Umbenennung, Feldparität und die masselose konstante Verschiebung.
Unter ausreichender glatter Kontrolle ist er im führenden Grenzfall
unterdrückt, verändert aber die endlichen Matrizen. Hier wird ausdrücklich
$\zeta=0$ gewählt. Eine eindeutige mikroskopische Wirkung oder eine unter
Quantenkorrekturen geschlossene Kopplungsfamilie ist damit nicht bewiesen.

\newpage
\sect{7. Evidenz, Abbruchgrenzen und nächster Anschluss}
\textbf{25 automatisierte Kontrollen.}
Die Prüfungen verwenden exakte rationale beziehungsweise symbolische
Algebra. Sie decken Produkt- und Paarform, Hesse-Matrix und erste
Variation, Symmetrien, Einheiten-relevante Normierung, Randreaktion,
Summationsidentität, gemeinsame Ordnungsvariation, Schnittbeiträge,
klassische Zielidentitäten und die kleinen Kontrollfamilien ab.
Über alle 107 deklarierten gerichteten Carrier werden die Symmetrie,
Konstantenreaktion, Green-Identität und Adapterreihenfolge kontrolliert.
Die Anzahl bestandener Tests ist kein physikalischer Konvergenznachweis.

Das Paket enthält den exakten Kandidaten in \texttt{model.py},
\texttt{action\_contract.json}, die Tests, maschinenlesbare Ergebnisse,
eine reproduzierbare Archivprüfung sowie unveränderte, gehashte Eingaben
des bisherigen roten Fadens. Grundlage ist der zusammengeführte Stand
\href{https://github.com/thomasmgfeller/x6-spectral-astrolabium/pull/278}{PR 278},
Commit\\{\small\texttt{4fed3a2fd0c5050af8e5b6dab6e6311fe7f4fd83}}.
Diese neue Ausarbeitung ändert diesen Commit nicht.

\textbf{Was Schritt 1 abschließt.}
Feldinhalt, minimale Kopplungswahl, Einheiten, endliche gemeinsame Wirkung,
skalare Randvariation und klassische Zielgleichungen sind jetzt explizit.
Für den gewählten Kandidaten existiert eine aus der Wirkung berechnete
Feldmatrix; ihre Nichtsingularität im Inneren der 33 AA-Kontrollen ist
bewiesen. Es wurden keine kostenpflichtigen oder großen Produktionsläufe
ausgeführt und keine Zahlen für physikalische Übergangseinträge eingesetzt.

\textbf{Was offen bleibt.}
Der Kandidat besitzt noch kein hergeleitetes physikalisches Historiemaß,
keine vollständige gravitative Randwirkung, keine qualifizierten
Randzustände und kein normiertes gemeinsames Quantenmaß. Insbesondere
fehlen die gemeinsame Kontinuums- und Variationskonvergenz,
Stabilitätskontrolle, Renormierungsfestlegung und Gesamtfehlerbilanz.
Deshalb bleiben \texttt{MODEL}, \texttt{CURRENT}, \texttt{CORRESPONDENCE},
\texttt{BOUNCE}, \texttt{REGULARITY}, \texttt{BOUNDARY}, \texttt{TRAPPING}
und \texttt{EXTENSION} sämtlich auf \status{HOLD}.

\textbf{Unmittelbarer Anschluss an Schritt 2.}
Jetzt können physikalische Randinformation und Zustände für diesen
\emph{benannten Kandidaten} spezifiziert werden: äußere Kragen,
Strata-Projektoren, Feldpakete und Skalarprodukte einschließlich ihrer
Maß- und Einheitenkonventionen. Bis dahin bleiben die physikalische
$8\times8$-Matrix und insbesondere $K_{0+}$ nicht konstruiert.
Die Wahl zusätzlicher Materie muss dabei sichtbar bleiben; sie ist keine
nachträgliche Umdeutung des bisherigen reinen Staubmodells.

\sect{Primärquellen und Reichweite ihrer Verwendung}
\small
[1] D. M. T. Benincasa, F. Dowker: \emph{The Scalar Curvature of a
Causal Set}, Phys. Rev. Lett. 104, 181301 (2010), korrigierte Fassung v4.
Verwendung: BD-Operator und Krümmungsbezug, kein Beweis dieses Gesamtmodells.\\
\smallurl{https://arxiv.org/abs/1001.2725}

[2] A. Belenchia, D. M. T. Benincasa, F. Dowker: \emph{The continuum
limit of a 4-dimensional causal set scalar d'Alembertian},
Class. Quantum Grav. 33, 245018 (2016).
Verwendung: mittlerer Operatorgrenzfall und dessen Voraussetzungen.\\
\smallurl{https://arxiv.org/abs/1510.04656}

[3] V. Husain, T. Pawłowski: \emph{Time and a physical Hamiltonian
for quantum gravity}, Phys. Rev. Lett. 108, 141301 (2012), Gleichung (1).
Verwendung: klassischer wirbelfreier Staub plus unabhängige Materie.\\
\smallurl{https://arxiv.org/abs/1108.1145}

[4] S. Aslanbeigi, M. Saravani, R. D. Sorkin: \emph{Generalized causal
set d'Alembertians}, JHEP 06 (2014) 024.
Verwendung: Operatorfamilien und Grenzen der Stabilitätsannahme.\\
\smallurl{https://arxiv.org/abs/1403.1622}
\end{document}
